\section{Introduction}
A posterior distribution contains more information than is needed for a
finite future experiment. The amount that can be discarded depends on
which posterior distributions the past can produce, which functions the
future can test, and whether a compressed state can be updated after the
next report. At a degeneracy these three questions need not have the
same answer. An observation matrix may lose a direction which the past
has scarcely acquired; a regular local image need not control all
histories; and a sequence of checkpoint encoders need not be a causal
filter.

We determine their joint effect for positive sparse monomial experiments.
The main invariant is a finite list of maximal Vandermonde products of
future exponents, truncated by the dimension acquired by the past. It
gives the sharp prediction error for every integer memory budget,
uniformly through all additive collisions, and the same law is achieved
by one causal filter. Thus the result is not merely an exponent at a
fixed nonsingular parameter: it determines the transition between
memory regimes as directions become indistinguishable.

The proof proceeds by a transfer principle formulated for reachable
sets, actual acquisition measures, and state updates. Its geometric
hypotheses are a dimension-truncated whole-image cover and positive
mass in every initial resolution flag. A uniform update bound then
turns all checkpoint covers into a single causal realization, with an
explicit error recurrence for unequal stage budgets. The geometric
verification for monomial experiments is the main model-specific
argument. It pairs a product tangent generated by real commands with
complete Newton--Hermite prefixes and retains the probability of the
report word throughout.

\subsection{The main theorem}
Let
\[
 A=\{0=a_0<a_1<\cdots<a_{r-1}\},\qquad
 W_A=\operatorname{span}\{t^{a_i}:0\le i<r\},\qquad r\ge2.
\]
The latent variable lies in $[0,1]$ and has a fixed full-support
probability $\mu$, not necessarily admitting a density. A fixed matrix
$C=(c_{ji})$ of column rank $r$ specifies a finite detector
\[
 k_j^a(t)=\sum_i c_{ji}t^{a_i},\qquad
 \sum_j k_j^a(t)=1,\qquad k_j^a(t)\ge k_*>0.
\]
The known exponent vector $a$ ranges over a compact subset $K$ of the
strictly ordered one-step chamber, and these inequalities hold on $K$.
A command $g\in[\eta,1-\eta]^J$, $0<\eta<1/2$, retains cell $j$
with probability $g_j$. The reports are the retained cell or failure;
all trials, including failures, are counted. Every report likelihood is
bounded below by a fixed $\kappa>0$. Acquisition commands are independent
and uniform for the unconditional criterion.

Fix a total horizon $N\ge2$ and choose
$H\ge\max\{1,N\max_K a_{r-1}\}$. At checkpoint $n$ there are
$m=N-n$ future trials. Keep every formal label
\[
 \Lambda_m^+=\{\alpha\in\N_0^r:|\alpha|=m\}\setminus\{me_0\},
 \qquad q_m=\binom{m+r-1}{r-1}-1,
 \qquad x_\alpha(a)=H^{-1}\alpha\cdot a,
\]
even when several labels have equal values. For $1\le\ell\le q_m$ put
\[
 \mathcal V_{m,\ell}(a)=\max_{\substack{F\subset\Lambda_m^+\\|F|=\ell}}
     \prod_{\{\alpha,\beta\}\subset F}|x_\alpha(a)-x_\beta(a)|,
 \qquad \mathcal V_{m,1}=1.
\]
A zero product is retained. With $p_{n,m}=\min\{n(r-1),q_m\}$ define
\[
 \Psi_{n,m}(M,a)=\max_{1\le\ell\le p_{n,m}}
       (\mathcal V_{m,\ell}(a)/M)^{2/\ell},\qquad
 \Xi_N(M,a)=\max_{1\le n<N}\Psi_{n,N-n}(M,a).
\]
Endpoint profiles are zero. The future query is sampled independently
\emph{after} the memory label is formed, from a fixed finite spanning
menu of actual $m$-trial products. Only that query is executed. Regret
is excess squared prediction loss for its binary outcome. A checkpoint
encoder can read the complete prefix before keeping one of $M$ labels;
a causal filter updates only its retained label using the current
command and report. Definition~\ref{def:finite-state} gives the precise
resource and loss conventions.

\begin{theorem}[Intrinsic attainable resolution and causal memory]\label{thm:resolution-main}
Under these hypotheses there are constants $0<c\le C<\infty$,
depending only on the fixed experiment, $K$, $N$ and $\mu$, with the
following properties for every known calibration $a\in K$ and every
integer $M\ge1$.

At each checkpoint $n+m=N$, the optimal $M$-message squared-prediction
regret is between $c\Psi_{n,m}(M,a)$ and $C\Psi_{n,m}(M,a)$,
both unconditionally under independent uniform acquisition commands and
in the minimax sense over admitted histories. The query is sampled
independently after the index is formed, from the ordered products of
one fixed attainable failure-factor basis; all its $m$ trials are counted.

There exists one deterministic stage-compatible transducer, allowed to
depend on $a$ and $M$, storing at most $M$ persistent labels after
every report, whose regret at every checkpoint and on every history is
at most $C\Xi_N(M,a)$. Conversely every such transducer has maximum
unconditional checkpoint regret at least $c\Xi_N(M,a)$ under the
same common exploration law. Independent coding randomization does not
remove the lower bound. The clock and fixed real-valued program are
read-only as in Definition~\ref{def:finite-state}.

The constants are uniform across all additive collision strata of $K$.
No calibration-blind codebook or uniformity over all full-support priors
is part of the assertion.
\end{theorem}


At a fixed calibration the largest nonzero index is
$\min\{n(r-1),|mA|-1\}$, the exact continuous-state dimension of
Theorem~\ref{thm:main}. Away from a collision this dimension determines
the small-error exponent. The full list $\mathcal V_{m,\ell}$ also
determines when that exponent becomes visible. In particular,
Corollary~\ref{cor:intrinsic-bits} recovers the persistent-bit law, and
Theorem~\ref{thm:collision-tree} gives the phases along collision paths
of arbitrary finite contact orders. Intersecting collision loci need
not reduce to one gap; the two-parameter example in
Section~\ref{sec:collision-geometry} gives distinct intermediate regimes.

\subsection{The implication from geometry to causality}
The transfer theorem, Theorem~\ref{thm:causal-transfer}, is stated for
an arbitrary family of compact reachable states with linear physical
query coordinates. After a uniformly invertible coordinate change,
suppose the entire state set has bounded semialgebraic format,
dimension at most $p_n$, and containing side lengths
$s_{n,1}\ge\cdots\ge s_{n,q_n}\ge0$. Suppose also that one common
exploration experiment supplies positive unconditional mass on a
rectangle in every initial nonzero flag of length at most $p_n$.
If the exact updates are uniformly Lipschitz and preserve reachability,
then the checkpoint error is
\[
 Q_n(M)=\max_{\ell\le p_n}
       \left(\frac{s_{n,1}\cdots s_{n,\ell}}{M}\right)^{2/\ell}.
\]
For arbitrary stage budgets $M_n$, one causal filter has error at most
a fixed-horizon multiple of $\max_n Q_n(M_n)$, and every filter has
error at least another such multiple under that same exploration law.
The theorem gives the full convolution of the preceding quantization
errors, rather than replacing the accumulated error by the last one.

This implication is useful precisely when well-conditioned geometric
coordinates have no well-conditioned recursion. In the monomial
problem, divided differences describe the resolution flags, but raw
remaining moments carry the filter. Positivity of a report likelihood
bounds the Bayes denominator along every segment between reachable
moment states. Lemma~\ref{lem:positive-remaining-update} supplies the
update bound without a reciprocal collision gap, on arbitrary compact
latent spaces with finite remaining-test closure. It does not assert
that every such experiment has the same attainable geometry.

The model-specific verification has three steps. First, the binomial
product tangent in Lemma~\ref{lem:binomial-tangent} has dimension
$n(r-1)+1$. Pairing it with a complete confluent prefix and removing
evidence loses exactly one direction. The inverse chart includes
kernel coordinates; integrating them while retaining the all-failure
probability proves Lemma~\ref{lem:newton-attainment}. Second, each
report-word image is rational in its command variables. Normalizing
each factor gives dimension at most $n(r-1)$, and the
thin-rectangle estimate of Lemma~\ref{lem:tame-rectangle} covers its
\emph{whole} image, not a section. Third, the remaining-test update
lemma supplies the causal hypothesis. Section~\ref{sec:transfer-application}
checks these hypotheses explicitly, including a bounded invertible
extension of the Newton coordinate change at exact collisions, and
proves Theorem~\ref{thm:resolution-main}.

The elementary cover-to-filter argument is not the source of a new
entropy inequality or a new interpolation identity. The substantive
compatibility is that all acquired flags needed for the lower bound,
the dimensional truncation of the global upper bound, and one stable
recursion remain valid at the same degeneracies. Without the acquired
measure one obtains only an observation-side estimate; without the
global cover one controls only selected histories; without the update
one obtains checkpoint encoders rather than a filter.

\subsection{Companion results and the organization of the proofs}
The principal proof is contained in Sections~\ref{sec:experiment}--\ref{sec:transfer-application},
with the bit and collision-path consequences immediately following it.
The appendices are organized by mathematical dependence and contain the
complete companion arguments, not just statements or proof sketches.
Appendix~\ref{app:structural} classifies the fixed-prior dimension
exponents of command-polynomial experiments and proves the uniform
one-step covariance profile. Appendix~\ref{app:pairings} treats
full-support moment pairings: a positivity criterion gives annihilation,
whereas the rank of a quotient multiplication pencil characterizes an
\emph{exact} kernel. Its square moving-kernel example separates this
criterion from feasibility, closure and dimension counting. The finite
unconstrained maximum-rank case is the classical linear
matroid-intersection calculation; the constrained compact-space result
uses an open chart of bounded density perturbations instead.

Appendix~\ref{app:additional-geometry} contains the circular-contrast
law, the operational inverse of the budget profile, and the joint
memory--prior-ambiguity geometry. Appendix~\ref{app:complete-routes}
contains the complete direct proofs of
Theorems~\ref{thm:intrinsic-checkpoint} and
\ref{thm:intrinsic-streaming}, together with their specialized analytic routes.
Appendix~\ref{app:decisions} records the observation algebra and decision
consequences. Appendix~\ref{app:effective} supplies the finite numerical
realization, resource and stability results. Precise comparisons with
classical inputs and neighboring approximation problems are collected
in Appendix~\ref{app:comparison}.

The scope distinctions matter. A rank formula for a specified prior
is not a multiscale formula uniform through rank loss. A local
exact-kernel theorem is not a global causal classification. The
different geometric classes and the uncertainty statements retain their
own assumptions, summarized in Table~\ref{tab:scope}. In particular,
the main memory resource counts persistent labels, not every bit of
the fixed program or arithmetic workspace. Its constants are uniform
in the calibration and integer budget at the stated fixed horizon,
not uniformly over all full-support priors or unbounded horizons.

\begin{table}[tbp]
\centering\footnotesize
\setlength{\tabcolsep}{4pt}
\begin{tabular}{@{}p{.16\textwidth}p{.23\textwidth}p{.34\textwidth}p{.21\textwidth}@{}}
\toprule
Result & Acquisition and latent space & Uniformity & Memory conclusion\\
\midrule
Collision law & Monomial, multiple steps; interval & Fixed full-support
prior; fixed horizon and compact ordered chamber; all additive collisions
& All integer budgets; checkpoint and one causal filter\\[3pt]
Structural rank & Command-polynomial; arbitrary compact space & Every
specified prior; uniform on compact constant-rank moment sets & Rank
exponent; checkpoint and one causal filter\\[3pt]
Affine profile & One affine step; arbitrary compact space & Across all
covariance rank losses, for fixed physical bounds and dimensions & Full
checkpoint profile; not a general multi-step profile\\[3pt]
Exact kernels & Continuous finite product spaces & Full-support
annihilation slice; arbitrarily small bounded positive tilts & Exact local
rank and acquired mass; global and causal claims require more\\[3pt]
Circular law & Multiple steps; circle and Haar prior & Fixed horizon and
known sufficiently small contrast, including zero & Paired scales;
checkpoint and one causal filter\\[3pt]
Prior ambiguity & Monomial moment uncertainty & Stated domination and
interior-slack hypotheses & Widths; joint memory and common-moment error\\
\bottomrule
\end{tabular}
\caption{Hypotheses and conclusions. In every memory statement the query
is chosen after the label is formed. Persistent labels are charged;
known calibration, clock, and the fixed program follow the stated
read-only convention. No discarded command tape is free.}
\label{tab:scope}
\end{table}
